Antes de condenar a Anderson, tres matices.

Primero, lo que espera en la caja es solo el bloque de control, la libreta del contador. El objeto en sí muere en el acto: su pico de objetos vivos es 14, como el de los demás.

Segundo, en bytes es poco: menos de medio mega.

Tercero, por eso la memoria total del proceso no lo delata: con un millón de pilas, lo que pesa son los datos.

Sin Docker sale lo mismo. Así que la fila 13 se reproduce: Anderson tiene un sobrecoste pequeño, y Autores, ninguno. Anderson compra velocidad de lectura con una caja de pendientes; Autores no la necesita.
